\documentclass[10pt,a4paper]{amsart}
\usepackage[margin=19mm]{geometry}
\usepackage{amsmath,amssymb,mathtools}
\usepackage[scr=rsfs,cal=euler]{mathalfa}
\usepackage{xfrac}
\usepackage[unicode,hidelinks,bookmarks=false]{hyperref}
\newcommand{\ie}{i.e.\ }
\newcommand{\cf}{cf.\ }
\newcommand{\resp}{respectively\ }
\newcommand{\Subsection}{\subsection}
\providecommand{\nobreakdash}{\nobreak\hbox{-}}
\setlength{\emergencystretch}{2em}
\multlinegap0pt
\usepackage{bm}
\usepackage{bbm}
\usepackage{dsfont}
\usepackage{xspace}
\usepackage{cancel}
\usepackage{mathtools}
\usepackage[shortlabels]{enumitem}
\usepackage{amsthm}
\makeatletter
\@ifundefined{cor}{\newtheorem{cor}{Corollary}}{}
\@ifundefined{prop}{\newtheorem{prop}{Proposition}}{}
\@ifundefined{thm}{\newtheorem{thm}{Theorem}}{}
\makeatother
\title[Clarke Differentials and the Envelope Theorem in Dynamic Pro]{Clarke Differentials and the Envelope Theorem in Dynamic Programming}
\author{Yuhki Hosoya}
\date{Source: arXiv:2509.17103v4 (CC BY 4.0)}
\begin{document}
\maketitle

\noindent\emph{Source and licence.} Clarke Differentials and the Envelope Theorem in Dynamic Programming. Version arXiv:2509.17103v4 (CC BY 4.0), distributed under the Creative Commons Attribution 4.0 International licence, as recorded in the arXiv licence metadata for this exact version and shown on the abstract page https://arxiv.org/abs/2509.17103v4. Full rights notice: licenses/arxiv-2509.17103-CC-BY-4.0.txt.\par\smallskip

\noindent\textit{Editorial scope.} Counted unit: the Cor whose source label is \texttt{cor3} in the article (source0.tex lines 444--447, source label \texttt{cor3}), delivered together with the article's own complete original proof of it (source0.tex lines 449--467, one \texttt{proof} environment of 19 lines). No mathematics is written, repaired or re-derived: statement and proof are the article's own text line for line. Required context retained uncounted: prop5 (lines 229--240); B (lines 307--309); thm3 (lines 353--366). Every definition, display and label of those lines is retained unchanged. Omitted from this package and not redistributed: 1--228, 241--306, 310--352, 367--443, 448--448, 468--784. Nothing inside a retained excerpt is omitted or altered: every retained body is delivered exactly as the article's lines are written, so no omission receipt of any kind accompanies this package. Citations: the retained excerpts carry no citation command at all, so no bibliography entry of the article is transcribed and no reference is redistributed. Reference adaptation: none was needed and none was made; every reference inside the delivered unit resolves inside it, so no unresolved reference or citation is produced. Printed slips kept literal: no printed word, symbol, hypothesis or number of the counted statement or of its proof was altered. Layout and class: the article's own class is \texttt{article} with packages including amsmath, amssymb, bm, graphicx, xcolor, enumerate, tikz, mathrsfs, amsthm; the wrapper is the lane's pinned \texttt{amsart} editorial wrapper at 10pt on A4, which restates the theorem-like environments the retained text needs and adds the article's own macro definitions that the retained text needs (none), with extra wrapper packages (bm, bbm, dsfont, xspace, cancel, mathtools, enumitem). The delivered numbering is the unit's own. Assets: no retained span contains a figure environment, a graphics-inclusion command, a PDF-page inclusion, a table or a TikZ picture; no image file is copied into this package, the package declares no asset dependency and no image file is redistributed with it. Licence: version arXiv:2509.17103v4 of this article is distributed under the Creative Commons Attribution 4.0 International licence, as recorded in the arXiv licence metadata for this exact version and shown on the abstract page https://arxiv.org/abs/2509.17103v4. A full rights notice is delivered at licenses/arxiv-2509.17103-CC-BY-4.0.txt. Characters: every retained excerpt is pure ASCII, so the delivered engine, XeLaTeX with its default Latin Modern text font, reports no missing character.
\begin{prop}\label{prop5}
Suppose that $U$ is a nonempty open subset of a Banach space $X$, $f:U\to \mathbb{R}$ is locally Lipschitz, and $x\in U$. Then, the following results hold.
\begin{enumerate}[(i)]
\item If $f$ is G\^ateaux differentiable at $x$, then $D_Gf(x)\in \partial^{\circ}f(x)$.

\item $\partial^{\circ}f(x)$ is a singleton if and only if $f$ is strongly differentiable. In this case, $\partial^{\circ}f(x)=\{D_sf(x)\}$.

\item If $f$ is strongly differentiable at $x$, then it is Hadamard differentiable at $x$, and $D_sf(x)=D_Hf(x)$.\footnote{If $X=\mathbb{R}^n$, then $D_Hf(x)=Df(x)$, and thus, $D_sf(x)=Df(x)$.}

\item If $f$ is strongly differentiable at $x$, then $f$ is regular at $x$.
\end{enumerate}
\end{prop}

\begin{equation}\label{B}
V(x)=\max\{w(x,y)+\delta V(y)|y\in \Gamma(x)\}.
\end{equation}

\begin{thm}\label{thm3}
Let $X$ be a Banach space, $W\subset X$ be nonempty and open, and $\Gamma:\bar{W}\twoheadrightarrow \bar{W}$ be nonempty-valued. Suppose that $w:T\to \mathbb{R}\cup \{-\infty\}$ is defined and continuous (where, $T=\{(x,y)\in \bar{W}^2|y\in \Gamma(x)\}$), and that $\bar{V}:\bar{W}\to \mathbb{R}\cup \{-\infty,+\infty\}$ is real-valued on $W$. Fix $\bar{x}\in W$, and suppose that the following requirements hold.
\begin{enumerate}[i)]
\item There exists an open neighborhood $W'\subset W$ of $\bar{x}$ such that $\Gamma$ is compact-valued and continuous on $W'$, and if $x,x'\in W'$, then $G(x)\cap \Gamma(x')\neq \emptyset$.

\item $w$ is real-valued and locally Lipschitz around $(x,y)$ for all $x\in W'$ and $y\in \mbox{cl.}G(x)$.

\item For every $\bar{y}\in \mbox{cl.}G(\bar{x})\cap W$, $w$ is regular at $(\bar{x},\bar{y})$.

\item $G(\bar{x})\subset W$ and $G$ is upper hemi-continuous at $\bar{x}$.\footnote{This is an additional requirement in this theorem. Because $w$ is not bounded, the Bellman operator is not a contraction in this theorem, and thus there may be a solution to (\ref{B}) other than $\bar{V}$. Hence, the usual arguments for verifying the continuity of $\bar{V}$ using the fixed point theorem cannot be applied in this theorem, and this additional requirement is needed. Note that, if $\bar{V}$ is continuous on some neighborhood of $G(\bar{x})$, then this requirement automatically holds.}
\end{enumerate}
Then, $\bar{V}$ is Lipschitz around $\bar{x}$, and for some $\bar{y}\in \mbox{cl.}G(\bar{x})$,
\[\partial^{\circ}\bar{V}(\bar{x})\subset \partial^{\circ}_xw(\bar{x},\bar{y}).\]
\end{thm}

\begin{cor}\label{cor3}
In addition to the assumptions of Theorem \ref{thm3}, suppose that $w_{\bar{y}}$ is strongly differentiable at $\bar{x}$, where $\bar{y}$ is given in Theorem \ref{thm3}. Then, $\bar{V}$ is also strongly differentiable and Hadamard differentiable at $\bar{x}$. Moreover, if $G(\bar{x})\subset \Gamma(x)$ for all $x\in W'$, then for any $y\in G(\bar{x})\cap W$, 
\[D_H\bar{V}(\bar{x})\in \partial^{\circ}_xw(\bar{x},y).\]
\end{cor}

\begin{proof}
By Proposition \ref{prop5}, $\partial^{\circ}_xw(\bar{x},\bar{y})$ is a singleton $\{D_sw_{\bar{y}}(\bar{x})\}$. Therefore,
\[\partial^{\circ}\bar{V}(\bar{x})=\{D_sw_{\bar{y}}(\bar{x})\}.\]
This implies that $\bar{V}$ is both strong differentiable and Hadamard differentiable at $\bar{x}$. Next, suppose that $G(\bar{x})\subset \Gamma(x)$ for all $x\in W'$. Choose any $y\in G(\bar{x})$ and $v\in X$. Let $(t_n)$ be a positive monotone sequence that converges to $0$. We can assume without loss of generality that $\bar{x}+t_nv\in W'$ for all $n$. Thus,
\[\bar{V}(\bar{x})=w(\bar{x},y)+\delta \bar{V}(y),\]
\[\bar{V}(\bar{x}+t_nv)\ge w(\bar{x}+t_nv,y)+\delta \bar{V}(y).\]
Hence,
\begin{align*}
\langle D_H\bar{V}(\bar{x}),v\rangle=&~-\lim_{n\to \infty}\frac{\bar{V}(\bar{x})-\bar{V}(\bar{x}+t_nv)}{t_n}\\
\le&~-\limsup_{n\to \infty}\frac{w(\bar{x},y)-w(\bar{x}+t_nv,y)}{t_n}\\
=&~\liminf_{n\to \infty}\frac{w(\bar{x}+t_nv,y)-w(\bar{x},y)}{t_n}\\
=&~w'(\bar{x},y;v,0)=w^{\circ}(\bar{x},y;v,0).
\end{align*}
By the same arguments as in the proof of Theorem \ref{thm3}, we have that
\[w^{\circ}(\bar{x},y;v,0)=w_y^{\circ}(\bar{x};v).\]
Therefore,
\[D_H\bar{V}(\bar{x})\in \partial^{\circ}_xw(\bar{x},y),\]
as desired. This completes the proof. $\blacksquare$
\end{proof}

\end{document}
